\section{Aplicación de redes neuronales: cuantificación de la pérdida}

\subsection{Problema de ejemplo: ¿aprobará el examen?}

Amini usa un ejemplo intuitivo para ilustrar el flujo de aplicación de las redes neuronales. El problema es: a partir del número de asistencias del estudiante $x_1$ y las horas dedicadas al proyecto final $x_2$, predecir si el estudiante aprobará el curso 6.S191.

\begin{figure}[H]
\centering
\includegraphics[width=0.85\textwidth]{slides-images/slide-047.jpg}
\caption{Problema de ejemplo: dado el nuevo punto de datos $[4,5]$, ¿se predice que aprobará?\protect\footnotemark}
\end{figure}
\footnotetext{Diapositivas, página 47.}

Al introducir el punto de datos $x^{(1)} = [4, 5]$ en una red con una sola capa oculta, supongamos que la red (con pesos inicializados aleatoriamente) produce $\hat{y}_1 = 0.1$. Pero el estudiante en realidad aprobó el curso (etiqueta real $y^{(1)} = 1$). Esto indica que los pesos actuales de la red aún no son suficientemente buenos: necesitamos una forma de cuantificar la diferencia entre la predicción y el valor real.

\subsection{Función de pérdida}

La \textbf{función de pérdida} (Loss Function) mide la diferencia entre el valor predicho por la red y el valor real.

Para la pérdida de un solo punto de datos:
$$\mathcal{L}\big(f(x^{(i)}; \mathbf{W}), y^{(i)}\big)$$

\begin{figure}[H]
\centering
\includegraphics[width=0.85\textwidth]{slides-images/slide-050.jpg}
\caption{Función de pérdida: mide la desviación entre la predicción y el valor real\protect\footnotemark}
\end{figure}
\footnotetext{Diapositivas, página 50.}

La \textbf{pérdida empírica} (Empirical Loss) sobre todo el conjunto de datos es el promedio de las pérdidas de todas las muestras:
$$J(\mathbf{W}) = \frac{1}{n}\sum_{i=1}^{n}\mathcal{L}\big(f(x^{(i)}; \mathbf{W}), y^{(i)}\big)$$

La pérdida empírica también se conoce como \textbf{función objetivo} (Objective Function), \textbf{función de costo} (Cost Function) o \textbf{riesgo empírico} (Empirical Risk).

\begin{importantbox}{Dos funciones de pérdida esenciales}
\textbf{1. Pérdida de entropía cruzada binaria} (Binary Cross Entropy Loss): se usa en problemas de clasificación (la salida es una probabilidad):
$$J(\mathbf{W}) = -\frac{1}{n}\sum_{i=1}^{n}\left[y^{(i)}\log f(x^{(i)};\mathbf{W}) + (1-y^{(i)})\log(1-f(x^{(i)};\mathbf{W}))\right]$$

\textbf{2. Pérdida de error cuadrático medio} (Mean Squared Error Loss): se usa en problemas de regresión (la salida es un valor continuo):
$$J(\mathbf{W}) = \frac{1}{n}\sum_{i=1}^{n}\left(y^{(i)} - f(x^{(i)};\mathbf{W})\right)^2$$
\end{importantbox}

\begin{figure}[H]
\centering
\begin{subfigure}[t]{0.48\textwidth}
\includegraphics[width=\textwidth]{slides-images/slide-052.jpg}
\caption{Pérdida de entropía cruzada (problema de clasificación)}
\end{subfigure}
\hfill
\begin{subfigure}[t]{0.48\textwidth}
\includegraphics[width=\textwidth]{slides-images/slide-053.jpg}
\caption{Error cuadrático medio (problema de regresión)}
\end{subfigure}
\caption{Dos funciones de pérdida esenciales\protect\footnotemark}
\end{figure}
\footnotetext{Diapositivas, páginas 52--53.}

\begin{warningbox}{Elegir la función de pérdida correcta}
La elección de la función de pérdida debe corresponder con el tipo de tarea:
\begin{itemize}
\item Las tareas de clasificación (aprobar/reprobar, gato/perro) usan la pérdida de entropía cruzada
\item Las tareas de regresión (predecir la calificación de un examen, predecir el precio de una casa) usan el error cuadrático medio
\end{itemize}
Una función de pérdida incorrecta impide que el modelo aprenda correctamente. Por ejemplo, usar el MSE en una tarea de clasificación provoca un entrenamiento lento porque la función de pérdida es poco sensible a los límites de probabilidad.
\end{warningbox}

\subsection{Resumen de la sección}

La función de pérdida es un componente central del entrenamiento de redes neuronales: cuantifica la diferencia entre la predicción del modelo y la etiqueta real. La pérdida empírica es el promedio de las pérdidas sobre todas las muestras de entrenamiento. Las tareas de clasificación suelen usar la pérdida de entropía cruzada, y las tareas de regresión, el error cuadrático medio. El objetivo del entrenamiento es encontrar los parámetros de peso que minimicen la pérdida empírica.
